\documentclass[11pt]{article}
\usepackage[letterpaper,margin=0.8in]{geometry}
\usepackage[T1]{fontenc}
\usepackage{mathptmx}
\usepackage{tabularx,booktabs,array}
\usepackage{enumitem}
\usepackage[table]{xcolor}
\usepackage{titlesec}
\usepackage[hidelinks]{hyperref}
\usepackage{fancyhdr}
\usepackage{needspace}
\definecolor{cmured}{HTML}{C41230}
\definecolor{rowgray}{HTML}{F3F3F3}
\titleformat{\section}{\large\bfseries\color{cmured}}{\thesection.}{0.5em}{}
\titleformat{\subsection}{\normalsize\bfseries}{\thesubsection}{0.5em}{}
\titlespacing*{\section}{0pt}{10pt}{3pt}
\titlespacing*{\subsection}{0pt}{6pt}{2pt}
\setlist{itemsep=1pt,topsep=2pt,leftmargin=1.3em}
\setlength{\parindent}{0pt}
\setlength{\parskip}{4pt}
\renewcommand{\arraystretch}{1.25}
\newcolumntype{L}{>{\raggedright\arraybackslash}X}
\newcolumntype{P}[1]{>{\raggedright\arraybackslash}p{#1}}
\pagestyle{fancy}
\fancyhf{}
\rhead{\small Group 19 \textbullet{} Intro to AI, Fall 2026}
\cfoot{\small\thepage}
\renewcommand{\headrulewidth}{0.4pt}

\usepackage{graphicx}
\hypersetup{pdftitle={Toolkit Activities: Step-by-Step Walkthrough},pdfauthor={Interactivity Foundation}}
\lhead{\small Toolkit Activities: Step-by-Step Walkthrough}
\rhead{\small Interactivity Foundation}
\sloppy
\newcounter{wstep}
\newcommand{\menu}[1]{\textbf{#1}}
\newcommand{\shot}[2]{\par\smallskip{\centering\fbox{\includegraphics[width=#2\textwidth,height=0.40\textheight,keepaspectratio]{walkthrough/#1}}\par}\smallskip}
\newcommand{\shots}[4]{\par\smallskip{\centering\fbox{\includegraphics[width=#3\textwidth,height=0.34\textheight,keepaspectratio]{walkthrough/#1}}\hfill\fbox{\includegraphics[width=#4\textwidth,height=0.34\textheight,keepaspectratio]{walkthrough/#2}}\par}\smallskip}
\newcommand{\say}[1]{\noindent\colorbox{rowgray}{\parbox{\dimexpr\textwidth-2\fboxsep}{\small\textbf{Say this:} \textit{#1}}}\par\medskip}
\newcommand{\wstep}[1]{\needspace{14\baselineskip}\stepcounter{wstep}\subsection*{Step \thewstep: #1}}
\begin{document}

\begin{center}
{\LARGE\bfseries Toolkit Activities: Step-by-Step Walkthrough}\\[3pt]
{\large Review, publish, add, and revise Toolkit activities on the IF website, one screen at a time}\\[2pt]
{\small Version 1, November 2026 \textbullet{} Prepared for IF staff by the CMU Heinz College student team \textbullet{} Companion to the Staff Training Guide}
\end{center}

\section*{How to use this walkthrough}
Each step says what to do, shows the screen with the thing to click outlined in red, and gives a short line to read aloud. Follow it on your own, or use it as the script for a short screen-recorded video.
\begin{itemize}
  \item \textbf{Three parts:} Part A, reviewing and publishing a copied draft (6 steps); Part B, adding a new activity (10 steps); Part C, revising a published activity and undoing a change (3 steps).
  \item \textbf{The golden rule} applies throughout: copy the author's words exactly. Never reword, shorten, or improve them.
  \item \textbf{You need} a WordPress login with the \textbf{Editor} or \textbf{Administrator} role.
\end{itemize}

\subsection*{Before you start: what these pictures show}
\begin{itemize}
  \item The pictures come from the team's test site: WordPress 6.8 with the Toolkit Activities add-on. Until IF's website manager installs the add-on on the IF website (pending as of November 2026), these screens exist only on a test site.
  \item The activities numbered 9.1, 9.2, and 9.3 are invented practice examples with made-up wording. There is no module 9, so they can't be mistaken for real Toolkit activities.
  \item The test site's design differs from the IF website. Check the words on the page, not the layout.
  \item In WordPress 6.5 and earlier, some names differ (for example, \menu{Save} is called \menu{Update}). The Staff Training Guide lists them.
\end{itemize}

\subsection*{Recording the videos}
\begin{itemize}
  \item Record on a test copy of the website, never on the live site, using only the practice activities. If IF has no test copy yet, wait until one is available.
  \item Make one short video per part (about 3 to 5 minutes each).
  \item \textbf{Mac:} press Shift, Command, and 5, choose \menu{Options > Microphone}, then \menu{Record Entire Screen}, then click \menu{Record}. \textbf{Windows 11:} open Snipping Tool, switch to video (record) mode, and turn on the microphone. \textbf{Windows 10, or if unsure:} record a Zoom or Teams meeting with your screen shared.
  \item Before recording: test your microphone, set the browser zoom to 100\%, close other tabs, hide the bookmarks bar, and turn off notifications so nothing private appears. Do one practice run.
  \item While recording: read the ``Say this'' line, move the mouse to the red-outlined item, then click. Pause for a second after each click, and leave two seconds of silence at the start and end for trimming.
\end{itemize}

\newpage
\section*{Part A: Review and publish a copied draft}
\setcounter{wstep}{0}
Use this part when the team has moved activities across and a draft is waiting for you.

\wstep{Find the draft}
In the left menu, click \menu{Toolkit Activities}. Drafts say \textbf{Draft} after the title. Click the draft's title (or hover over it and click \menu{Edit}).
\shot{A1_list.jpg}{0.78}
\say{In the menu on the left, I click Toolkit Activities. Every activity is listed here. This one says Draft, so it isn't on the public website yet. I'll click its title to open it.}

\wstep{Read the draft}
There is nothing to click in this step. The activity's text fills the middle of the screen, and its settings are on the right. In another browser tab, open the old Toolkit site (collaborativediscussionproject.com) and find the same activity, so you can compare them.
\shot{A2_editor.jpg}{0.78}
\say{This is the editor. The text is in the middle and the settings are on the right. I'll compare it with the same activity on the old Toolkit site.}

\wstep{Check the module, level, tags, and Activity details}
In the right-hand panel, scroll down to \menu{Modules}, \menu{Levels}, and \menu{Tags (formats)}. If a section is closed, click its name to open it. Check the right boxes are ticked. Below them, check \menu{Activity details}: Number, Time in minutes, Submitted by, PDF link, and Image credit.
\shots{A3_sidebar_1.jpg}{A3_sidebar_2.jpg}{0.30}{0.30}
\say{On the right, I check the module, the level, and the tags. Below them, in Activity details, I check the number, the time, who submitted it, the PDF link, and the image credit.}

\wstep{Preview the page}
At the top right, click the \menu{View} icon (it looks like a laptop screen). In the menu that opens, click \menu{Preview in new tab}.
\shots{A5a_view_icon.jpg}{A5b_view_menu.jpg}{0.47}{0.47}
A new tab shows the page as visitors will see it (the test site's design adds a line such as ``Written by in''; ignore it):
\shot{A5c_preview.jpg}{0.62}
\say{At the top right, I click the View icon, then Preview in new tab. This shows how the page will look to visitors.}

\wstep{Go through the checklist}
With the preview open, go through the review checklist in Section 6 of the Staff Training Guide, comparing with the old Toolkit page. If anything fails, don't publish; leave it as a draft and follow Section 3 of the guide.
\say{Now I go through the checklist in the Staff Training Guide. I'll pause the recording while I do.}

\wstep{Publish}
Back in the editor, click \menu{Publish} at the top right. A panel asks ``Are you ready to publish?''. Click \menu{Publish} again in that panel. (If no panel opens, the activity is already published.)
\shots{A6a_publish_button.jpg}{A6b_publish_panel.jpg}{0.47}{0.40}
When it's done, the panel says the activity \textbf{is now live}:
\shot{A6c_published.jpg}{0.36}
\say{Everything checks out, so I click Publish. WordPress asks if I'm ready, and I click Publish again. The panel says it's now live.}

\newpage
\section*{Part B: Add a new activity}
\setcounter{wstep}{0}
Before you start, have the activity's printable PDF ready, and prepare the text in the standard layout as a \textbf{plain-text} file: by hand with the fill-in template, or with Claude if IF has approved it (not yet approved as of November 2026). Wherever the activity has slides or a video, the prepared text has a line like \texttt{[Embed: link]}. See Section 4 of the Staff Training Guide.

\wstep{Upload the PDF}
In the left menu, go to \menu{Media}, then \menu{Add Media File}. Drag the PDF into the dotted box, or click \menu{Select Files}. Anyone with a file's link can open it once it is uploaded. If the upload fails because the file is too big, ask the website manager.
\shot{B1_media_add.jpg}{0.78}
\say{First, the printable PDF. I go to Media, then Add Media File, and drag the PDF into the box.}

\wstep{Copy the PDF's link}
When the upload finishes, click \menu{Copy URL to clipboard} next to the file. Keep this browser tab open: you'll copy the link again in Step 7.
\shot{B2_copy_url.jpg}{0.78}
\say{Once it's uploaded, I click Copy URL to clipboard. I'll keep this tab open, because I'll need the link again in a moment.}

\wstep{Start a new activity and add the title}
In another tab, go to \menu{Toolkit Activities > Add New Activity}. Click \textbf{Add title} and type the number, a space, and the title, for example \textit{9.3 Reframing Open Questions (practice)}.
\shot{B3_title.jpg}{0.78}
\say{Now I go to Toolkit Activities, Add New Activity. In the title box I type the number, a space, then the title.}

\needspace{22\baselineskip}\wstep{Paste the activity text}
Press Enter after the title, then paste everything from the description down to Additional Resources. Copy it from a plain-text file (Notepad, TextEdit in plain-text mode, or Claude's copy button), not from Word or Google Docs. The prepared text must have a blank line between every paragraph, heading, list, and \texttt{[Embed: link]} line; otherwise WordPress joins lines together. Lines starting with \texttt{\#} become headings and lines starting with \texttt{-} become bullet points. If you still see \texttt{\#} or \texttt{-} at the start of lines, undo (Ctrl+Z or Command+Z) and paste again from a plain-text copy. Do not change any words.
\shot{B4_pasted.jpg}{0.78}
\say{Below the title, I paste the activity text from a plain-text file. The headings and bullet points come in by themselves. I don't change a single word.}

\wstep{Add a YouTube or Vimeo video (if the activity has one)}
Click in the \texttt{[Embed: link]} line, select its text, and delete it, so the video goes in that spot. Then click the blue \menu{+} at the top left (it turns into an \textbf{X} while the block list is open), type \textit{YouTube} or \textit{Vimeo} in the search box, and click \menu{YouTube Embed} or \menu{Vimeo Embed}. Paste the video's normal watch link (for example youtube.com/watch?v=\ldots) and click \menu{Embed}.
\shot{B5a_embed_line.jpg}{0.62}
\shots{B5b_plus.jpg}{B5c_youtube.jpg}{0.47}{0.47}
\say{Where the activity has a video, I delete the Embed line, click the plus button, search for YouTube, and choose YouTube Embed. Then I paste the video's normal link.}

\wstep{Add Google Slides (if the activity has them)}
In Google Slides, choose \menu{File > Share > Publish to web}, open the \menu{Embed} tab, click \menu{Publish}, confirm with \menu{OK}, and copy the code. In WordPress, delete the \texttt{[Embed: link]} line as in Step 5, click \menu{+}, search \textit{Custom HTML}, choose it, and paste the code. For a practice run without a real slide deck, just delete the line.
\shot{B6_custom_html.jpg}{0.50}
\say{For Google Slides, I first publish the deck and copy its embed code in Google Slides. Then I click the plus button, choose Custom HTML, and paste the code in.}

\wstep{Fill in Activity details}
In the right-hand panel, scroll to \menu{Activity details}. Fill in Number, Time (numbers only, for example \textit{60}), and Submitted by. For the PDF link, go back to the Media tab from Step 2, click \menu{Copy URL to clipboard} again, and paste it here. Add the Image credit if there is one.
\shots{B7a_details_empty.jpg}{B7_details.jpg}{0.34}{0.34}
\say{In Activity details, I fill in the number, the time in minutes, and who submitted it. I copy the PDF link again from the Media tab and paste it in, and add the image credit if there is one.}

\wstep{Tick the module, level, and tags}
Find \menu{Modules}, \menu{Levels}, and \menu{Tags (formats)} in the right-hand panel (open a section by clicking its name if it's closed) and tick the right boxes. Use only the existing tags; don't use the \textbf{Add Category} links. A new tag needs an IF administrator's approval.
\shots{B8_ticks_1.jpg}{B8_ticks_2.jpg}{0.30}{0.30}
\say{Next I tick the module, the level, and the tags. I only use the tags that are already there.}

\wstep{Add the featured image (if the activity has a photo)}
Near the top of the right-hand panel, click \menu{Set featured image}, then upload the photo or choose it from the Media library.
\shot{B9_featured.jpg}{0.34}
\say{If the activity has a photo, I click Set featured image and choose it.}

\wstep{Preview, check, and publish}
{\raggedright Preview the page, go through the checklist, then click \menu{Publish} twice (see Part A, Steps 4 to 6).\par}
\say{Finally, I preview the page and go through the checklist. Then I click Publish, and Publish again.}

\newpage
\section*{Part C: Revise a published activity, and undo a change}
\setcounter{wstep}{0}

\wstep{Find the activity and click Edit}
Go to \menu{Toolkit Activities}, type the activity's number (for example \textit{9.2}) in the search box, and click \menu{Search Activities}. Hover over the title and click \menu{Edit}. Never use Add New Activity to change an existing activity; that creates a second page with a new address.
\shot{C1_search.jpg}{0.78}
\say{To change an activity that's already live, I search for its number, hover over it, and click Edit. I never make a new copy.}

\wstep{Make the change, preview, and save}
Change only the text that changed. If there is a new PDF, upload it (Part B, Steps 1 and 2) and replace the PDF link in Activity details. Preview it with the \menu{View} (laptop) icon, then \menu{Preview in new tab}. Don't use the arrow icon to its left: that opens the live page without your changes. Go through the checklist, then click \menu{Save}. The change is live straight away.
\shot{C2_save.jpg}{0.50}
\say{I change only what the author changed. I preview it with the laptop icon, check it, and click Save. On a published activity, the change goes live straight away.}

\wstep{Go back to an earlier version}
In the right-hand panel, click \menu{Revisions} (the number shows how many versions are saved). If you don't see Revisions, the activity hasn't been changed since it was published, so there is nothing to go back to.
\shot{C3_revisions_row.jpg}{0.36}
The Revisions screen shows what each version changed. The page's words appear with some code around them; ignore the code and look at the shaded lines: red (left) is what was removed, green (right) is what was added. Here, a step time was changed from 10 to 15 minutes:
\shot{C4a_diff.jpg}{0.66}
To undo that change, click \menu{Previous} once to move to the version before it (it greys out when there is no earlier version). The screen may then show that whole earlier version in green; that's normal. Click \menu{Restore This Revision}. On a published activity, this changes the live page at once. (Pending: the team confirms on the test site before December 3 that this also restores the Activity details box.)
\shot{C4b_restore.jpg}{0.78}
\say{If I need to undo a change, I click Revisions. This screen shows what changed: red was removed and green was added. I click Previous to go to the version before the mistake, then Restore This Revision.}

\section*{Need help?}
See the Staff Training Guide for the review checklist, what the review notes mean, common questions, and who to contact.

\end{document}
